\documentclass[11pt, a4paper]{article}
\usepackage[utf8]{inputenc}
\usepackage[ngerman]{babel}
\usepackage[T1]{fontenc}
\usepackage{amsmath, amssymb, amsfonts, amsthm}
\usepackage{geometry}
\geometry{a4paper, left=2.5cm, right=2.5cm, top=3cm, bottom=3cm}
\usepackage{hyperref}
\usepackage{titlesec}
\usepackage{enumitem}
\usepackage{setspace}
\usepackage{booktabs}
\setstretch{1.15}
\titleformat{\section}{\large\bfseries\sffamily}{\thesection}{1em}{}
\titleformat{\subsection}{\normalsize\bfseries\sffamily}{\thesubsection}{1em}{}
\newcommand{\mvec}[1]{\boldsymbol{#1}}
\newcommand{\clc}{\mathcal{C}\ell_{3}(\mathbb{C})}
\newcommand{\n}{\mvec{n}}
\newcommand{\EG}{\mathcal{E}_\Gamma}

\begin{document}
\begin{center}
  \sffamily
  {\Large\textbf{Clifford-BEM: Ausarbeitung zu AP 2}}\\[0.3cm]
  {\large Teil I: 3D-Galerkin-Prototyp, Kugel, Würfel, Vorkonditionierung}\\[0.3cm]
  {\small Arbeitsstand, Version 0.3}
\end{center}

\section*{Überblick}
Dieses Dokument behandelt die Schritte 2.1 und 2.2 des AP-2-Plans, den ersten Teil
von 2.3 sowie die Schritte 2.4 und 2.5 am Würfel (Abschnitte~\ref{sec:wuerfel}
und~\ref{sec:vorkond3d}). Gegenstand ist ein Python-Prototyp, der die Formulierung
$T_1 = E_2^+ + E_1^-J$ aus AP~1 in 3D mit dem Galerkin-Verfahren diskretisiert.
Die Skripte liegen in \texttt{ap2/}. Grundlage sind die Konventionen und die
Standardwahl der Hilfsparameter aus dem AP-1-Arbeitspapier.

\paragraph{Ergebnisse.} Der Randoperator trennt innere und äußere Spuren mit
Konvergenz erster Ordnung. Die Extinktion konvergiert für Glas, Gold und nahe der
Plasmonresonanz gegen Mie; nach Extrapolation beträgt die Abweichung $0{,}2$ bis
$0{,}7\,\%$. Der kleinste Singulärwert von $T_1$ bleibt bei Verfeinerung beschränkt. Die
GMRES-Iterationen sind für Glas konstant (12 bis 14) und wachsen für Gold nur schwach
(29 bis 37). Am Würfel bestätigt sich die Kantenvorhersage aus AP~1 in 3D quantitativ:
Für $\kappa=-3+0{,}3i$ fällt $\sigma_{\min}$ mit der Rate $0{,}31$, vorhergesagt ist
$0{,}32$. Gold und Glas bleiben stabil. Für $\kappa=-4+0{,}3i$ dominieren die Ecken. Mit einer Kanten- und
Eck-Blockvorkonditionierung sind die GMRES-Iterationen am Würfel unabhängig von der
Eckauflösung, für alle getesteten Kontraste (Abschnitt~\ref{sec:vorkond3d}).

\section{Diskretisierung (Schritt 2.1)}

\paragraph{Ansatzraum.} Ebene Dreiecke $\tau_i$ einer Ikosaeder-Kugel mit $n$-facher
Kantenunterteilung ($20n^2$ Dreiecke, Normalen nach außen). Die Dichten sind
stückweise konstant und $\clc$-wertig, mit der orthonormalen Basis
$\chi_{\tau_i}/\sqrt{|\tau_i|}$ und 8 Komponenten je Dreieck. Die Galerkin-Testung erfolgt
im $L^2$-Skalarprodukt $\langle A^\dagger B\rangle_0$; in dieser Basis ist die
Massenmatrix die Identität. $J$ wirkt dreiecksweise mit der Dreiecksnormalen und ist
blockdiagonal exakt darstellbar.

\paragraph{Galerkin-Einträge.} Für $\EG h=-2\,\mathrm{p.v.}\!\int\mvec\Phi_k(x-y)\n(y)h(y)\,dS_y$ ist
\[
  (E)_{ij} = -\frac{2}{\sqrt{|\tau_i||\tau_j|}}\;L\Big(\Big[\int_{\tau_i}\!\int_{\tau_j}\mvec\Phi_k(x-y)\,dS_y\,dS_x\Big]\,\n_j\Big)\ \in\mathbb{C}^{8\times8},
\]
mit der Linksmultiplikation $L(\cdot)$. Der Kern wird zerlegt in
\[
  \mvec\Phi_k(z) = \underbrace{\frac{z}{4\pi|z|^3}}_{\mvec\Phi_0}\ \underbrace{-\,\frac{ik}{4\pi|z|}}_{\text{schwach singulär}}\ +\ \mathrm{Rem}_k(z),
\]
wobei $\mathrm{Rem}_k$ beschränkt ist (AP~1, Lemma~3.1).

\paragraph{Behandlung der Singularitäten.}
\begin{itemize}[leftmargin=*]
  \item \emph{Fernpaare:} Gauß-Produktregel $7\times7$ (Dunavant, Grad 5) für den
    vollen Kern.
  \item \emph{Nahpaare} (Schwerpunktabstand unter $2{,}5\,h$): Das innere Integral über
    $\tau_j$ wird für $\mvec\Phi_0$ und $1/|z|$ analytisch berechnet. Für den
    Vektoranteil gilt
    \[
      \int_{\tau_j}\frac{x-y}{|x-y|^3}\,dS_y = \sum_{e}\mvec{m}_e f_e + \operatorname{sign}(w)\sum_e\beta_e\,\n_j,
      \qquad f_e=\ln\frac{R_e^++l_e^+}{R_e^-+l_e^-},
    \]
    mit äußeren Kantennormalen $\mvec{m}_e$, Höhe $w$ und den Winkeln $\beta_e$ der
    Wilton-Formeln. Die Normalkomponente ist der vorzeichenbehaftete Raumwinkel.
    Außen wird mit $4^2\cdot7=112$ Punkten integriert; $\mathrm{Rem}_k$ ist glatt genug für
    die Gauß-Regel. Kontrolle gegen feine Quadratur: relativer Fehler
    $10^{-15}$ bis $10^{-12}$ (\texttt{analytic.py}).
  \item \emph{Selbstterm:} Der $\mvec\Phi_0$-Anteil verschwindet exakt, weil
    $\int_\tau\int_\tau (x-y)/|x-y|^3$ antisymmetrisch ist. Der Hauptwert ist damit ohne
    Quadratur erledigt.
\end{itemize}
Die Assemblierung ist vektorisiert. Für $N=500$ Dreiecke (4000 Unbekannte) dauert sie
etwa 11\,s je Wellenzahl (\texttt{assemble.py}).

\section{Validierung (Schritt 2.2)}

\paragraph{Spurtrennung.} Für Spuren innerer bzw.\ äußerer Dirac-Lösungen
$\mvec\Phi_k(x-x_s)c$ (Quelle außen bzw.\ innen, $c$ zufällig) muss $Eh=\pm h$ gelten.
Auf dem Polyeder gilt die Cauchy-Formel exakt, der Test prüft also nur die
Diskretisierung (\texttt{test\_plemelj.py}):
\begin{center}\small
\begin{tabular}{lcccccc}
\toprule
& \multicolumn{2}{c}{$k=0$} & \multicolumn{2}{c}{$k=1{,}5$} & \multicolumn{2}{c}{$k=1{,}2+0{,}4i$}\\
$N$ & innen & außen & innen & außen & innen & außen\\
\midrule
80 & $4{,}6\cdot10^{-2}$ & $4{,}6\cdot10^{-2}$ & $5{,}0\cdot10^{-2}$ & $3{,}1\cdot10^{-2}$ & $5{,}4\cdot10^{-2}$ & $3{,}8\cdot10^{-2}$\\
180 & $2{,}3\cdot10^{-2}$ & $2{,}2\cdot10^{-2}$ & $2{,}5\cdot10^{-2}$ & $1{,}5\cdot10^{-2}$ & $2{,}6\cdot10^{-2}$ & $1{,}8\cdot10^{-2}$\\
320 & $1{,}4\cdot10^{-2}$ & $1{,}4\cdot10^{-2}$ & $1{,}5\cdot10^{-2}$ & $0{,}9\cdot10^{-2}$ & $1{,}6\cdot10^{-2}$ & $1{,}1\cdot10^{-2}$\\
\bottomrule
\end{tabular}
\end{center}
Das ist Konvergenz erster Ordnung in $h$, wie für stückweise konstante Ansätze zu
erwarten. Vorzeichen, Orientierung und Singularitätsbehandlung stimmen.

\paragraph{Streuung gegen Mie.} Gelöst wird $T_1h=h^{\mathrm{inc}}$ für eine ebene Welle,
mit dichtem LU. $Q_{\mathrm{ext}}$ folgt über das optische Theorem aus dem Fernfeld
\[
  F_\infty(\hat{\mvec x}) = -\frac{ik_2}{4\pi}\int_\Gamma e^{-ik_2\hat{\mvec x}\cdot y}(1+\hat{\mvec x})\,\n\,h^s\,dS
\]
(\texttt{scatter3d.py}). Die Extrapolation nimmt $Q(n)=Q_\infty-Cn^{-p}$ an und bestimmt
$p$ aus $n=3,4,5$:
\begin{center}\small
\begin{tabular}{lccccccc}
\toprule
Fall & $N=80$ & 180 & 320 & 500 & $p$ & extrapoliert & Mie\\
\midrule
Glas, $\omega a=1$ & 0,1707 & 0,1940 & 0,2029 & 0,2072 & 1,85 & 0,2156 & 0,2151\\
Gold, $\omega a=0{,}5$ & 0,8053 & 0,7107 & 0,6692 & 0,6473 & 1,51 & 0,5928 & 0,5900\\
$-2{,}2+0{,}3i$, $\omega a=0{,}3$ & 7,00 & 8,89 & 9,75 & 10,20 & 1,50 & 11,34 & 11,26\\
\bottomrule
\end{tabular}
\end{center}
Die Abweichung nach Extrapolation beträgt $0{,}24\,\%$, $0{,}47\,\%$ bzw.\ $0{,}73\,\%$. Der
größte Teil des Fehlers vor der Extrapolation ist geometrisch: Das eingeschriebene
Polyeder hat bei $N=80$ nur $93\,\%$ der Kugeloberfläche, und plasmonische Antworten
reagieren empfindlich auf die Form. Die beobachteten Ordnungen $p\approx1{,}5$ bis $1{,}9$
liegen zwischen der $L^2$-Ordnung~1 der Dichten und der Ordnung~2 des
Geometriefehlers.

\paragraph{Hilfskomponenten.} Ihr Anteil an der Lösung sinkt mit $h$: für Glas von
$2{,}8\cdot10^{-3}$ ($N=80$) auf $5{,}6\cdot10^{-4}$ ($N=500$), für Gold von $5{,}3\cdot10^{-3}$ auf
$1{,}1\cdot10^{-3}$. Das passt zur Eindeutigkeit mit Standardwahl (AP~1, Satz~10.2): Die
Hilfsanteile sind Diskretisierungsfehler und verschwinden im Grenzwert.

\section{Diskrete Stabilität auf der Kugel (Schritt 2.3, erster Teil)}
GMRES ohne Neustart bis $10^{-8}$, ohne und mit der punktweisen Vorkonditionierung
$2(1+J)^{-1}$, sowie der kleinste Singulärwert von $T_1$ in der orthonormalen Basis
(\texttt{stability3d.py}):
\begin{center}\small
\begin{tabular}{lcccc}
\toprule
$N$ & 80 & 180 & 320 & 500\\
\midrule
Glas: GMRES ohne / mit & 13 / 12 & 13 / 12 & 13 / 12 & 14 / 12\\
Glas: $\sigma_{\min}(T_1)$ & 0,457 & 0,450 & 0,448 & --\\
Gold: GMRES ohne / mit & 34 / 29 & 36 / 32 & 37 / 34 & 37 / 35\\
Gold: $\sigma_{\min}(T_1)$ & 0,223 & 0,216 & 0,213 & --\\
\bottomrule
\end{tabular}
\end{center}
Der kleinste Singulärwert ist in beiden Fällen gleichmäßig beschränkt, die
Galerkin-Diskretisierung ist auf der Kugel also stabil. Für Glas sind die
Iterationszahlen konstant. Für Gold wachsen sie schwach und scheinbar
sättigend. Die punktweise Vorkonditionierung spart nur wenige Iterationen, und ihr
Vorteil schrumpft bei Gold mit $N$. Eine plausible Ursache sind die Kanten des Polyeders:
Ihre Diederwinkel nähern sich mit $N$ zwar $\pi$, sind aber nicht glatt. Nach AP~1
liegt Gold für Winkel nahe $\pi$ weit außerhalb der kritischen Mengen, eine
Instabilität ist also nicht zu erwarten. Das Spektrum wird durch die Kanten aber
aufgefächert, wie im 2D-Prototyp.

\section{Würfel mit Kanten und Ecken (Schritt 2.4)}\label{sec:wuerfel}

\paragraph{Netz.} Würfel $[-1,1]^3$. Jede Seitenfläche trägt ein Tensornetz, das
geometrisch zu den Kanten hin gradiert ist. Die Knoten je Halbseite sind
$0;\ 0{,}5;\ 1-2^{-2};\dots;\ 1-2^{-L};\ 1$, der kleinste Kantenabstand ist also $2^{-L}$.
Die Quadrate werden so in Dreiecke geteilt, dass das Netz unter allen drei
Koordinatenspiegelungen invariant ist. Das ergibt $48(L+1)^2$ Dreiecke, für $L=5$ also
1728 Dreiecke und 13\,824 Unbekannte (\texttt{cube\_mesh.py}). Die Elemente an den Kanten
sind anisotrop.

\paragraph{Symmetriereduktion.} Die Spiegelungen $x_i\mapsto-x_i$ wirken auf
Multivektorfelder als $(U_ih)(x)=\mvec{e}_i\,\hat h(S_ix)\,\mvec{e}_i$, mit der
Hauptinvolution $\hat h$. Physikalisch transformiert der Vektoranteil ($\mvec{E}$) polar und
der Bivektoranteil ($I\mvec{H}$) axial. $E_k$ und $J$ kommutieren mit $U_i$; der Kommutator
ist $2\cdot10^{-15}$ relativ, bei $J$ exakt null. Die acht Charaktere der Gruppe
$\{\pm1\}^3$ zerlegen $T_1$ in acht Sektoren,
\[
  T_\chi[s,t]=\sum_{g}\chi(g)\,T[s,g\,t]\,A_g,\qquad s,t\in\mathcal{F},
\]
wobei $\mathcal{F}$ die Dreiecke im Oktanten $x,y,z>0$ sind und $A_g$ die $8\times8$-Wirkung
von $g$. Assembliert werden nur die Zeilen zu $\mathcal{F}$, also $1/8$ der Matrix. Für
$L=1$ stimmt $\min_\chi\sigma_{\min}(T_\chi)$ mit $\sigma_{\min}$ der vollen Matrix auf
alle Stellen überein (\texttt{cube\_sym.py}).

\paragraph{Nahfeld bei starker Gradierung.} Bei Paaren sehr ungleich großer Dreiecke
wird über das kleinere Dreieck außen integriert und über das größere innen analytisch
($\mvec\Phi_0$ ist ungerade, der skalare Anteil gerade). Für $L=5$ gibt es in den
assemblierten Zeilen 70\,570 Nahpaare; die Assemblierung beider Wellenzahlen dauert
etwa 65\,s.

\paragraph{Ergebnisse} ($\omega a=0{,}5$, $a=1$ halbe Kantenlänge; $\min$ über alle
Sektoren; gewichtet heißt Ähnlichkeit mit $r^{0{,}35}$, $r$ = Abstand zur nächsten
Kante; GMRES bis $10^{-8}$, ohne Vorkonditionierer, Maximum über die Sektoren;
\texttt{cube\_run.py}):
\begin{center}\small
\begin{tabular}{llccccc}
\toprule
$\kappa$ & & $L=1$ & 2 & 3 & 4 & 5\\
\midrule
Glas $2{,}25$ & $\sigma_{\min}$, $L^2$ & 0,584 & 0,582 & 0,581 & 0,581 & 0,580\\
 & GMRES & 11 & 11 & 11 & 11 & 12\\
\midrule
Gold $-11+1{,}2i$ & $\sigma_{\min}$, $L^2$ & 0,117 & 0,099 & 0,090 & 0,085 & 0,082\\
 & $\sigma_{\min}$, gewichtet & 0,118 & 0,102 & 0,094 & 0,090 & 0,087\\
 & GMRES & 43 & 48 & 50 & 54 & 56\\
\midrule
$-4+0{,}3i$ & $\sigma_{\min}$, $L^2$ & 0,087 & 0,061 & 0,048 & 0,040 & 0,033\\
 & $\sigma_{\min}$, gewichtet & 0,086 & 0,063 & 0,055 & 0,051 & 0,048\\
 & GMRES & 42 & 49 & 55 & 62 & 67\\
\midrule
$-3+0{,}3i$ & $\sigma_{\min}$, $L^2$ & 0,096 & 0,078 & 0,061 & 0,048 & 0,039\\
 & $\sigma_{\min}$, gewichtet & 0,096 & 0,082 & 0,069 & 0,058 & 0,051\\
 & GMRES & 41 & 50 & 58 & 65 & 72\\
\bottomrule
\end{tabular}
\end{center}

\paragraph{Vergleich mit den Vorhersagen aus AP~1.} Die $L^2$-Rate
$\sigma_{\min}\propto h_{\min}^{\gamma}$ ergibt sich an Kanten aus
$\gamma=\tfrac12-\operatorname{Re}a_{\min}$. An der Ecke ist als grober Vergleich der Kegel
mit gleichem Raumwinkel ($41{,}4^\circ$) herangezogen, mit $\gamma=-\operatorname{Re}\nu$:
\begin{center}\small
\begin{tabular}{lcccc}
\toprule
$\kappa$ & Kante: $\gamma$ & Vergleichskegel: $\gamma$ & beobachtet ($L=4\to5$) & Kante in $L^2$-Schleife?\\
\midrule
Glas & -- & -- & $0{,}001$ & nein\\
Gold & -- (Exp.\ $0{,}59>\tfrac12$) & -- (kein Exp.) & $0{,}06$, fallend & nein\\
$-4+0{,}3i$ & $0{,}12$ & $0{,}39$ & $0{,}28$ & ja\\
$-3+0{,}3i$ & $0{,}32$ & $0{,}34$ & $0{,}31$ & ja\\
\bottomrule
\end{tabular}
\end{center}
\begin{itemize}[leftmargin=*]
  \item \emph{Glas und Gold} sind stabil. Für Gold nimmt die Abnahme von Stufe zu Stufe ab
    (Quotienten $0{,}85$; $0{,}91$; $0{,}94$; $0{,}96$), verträglich mit einem positiven
    Grenzwert. Weder an der Kante noch am Vergleichskegel gibt es für Gold einen
    Singulärexponenten im kritischen Streifen.
  \item \emph{$-3+0{,}3i$:} Die beobachtete Rate $0{,}31$ stimmt mit der Kantenvorhersage
    $0{,}32$ überein (Kegel $0{,}34$). Das ist die erste quantitative Bestätigung der
    Kantentheorie in 3D.
  \item \emph{$-4+0{,}3i$:} Die Rate $0{,}28$ liegt deutlich über der Kantenvorhersage
    $0{,}12$, aber unter dem Vergleichskegel $0{,}39$. Die Ecken dominieren hier also. Der
    Vergleichskegel ist nur eine Näherung der echten Würfelecke, und bei $h_{\min}=2^{-5}$
    ist die Rechnung noch präasymptotisch.
  \item \emph{Gewichtet} ($r^{0{,}35}$ mit dem Kantenabstand): Für $-4+0{,}3i$ flacht die
    Abnahme deutlich ab (Quotienten $0{,}73$; $0{,}87$; $0{,}93$; $0{,}94$), der
    $L^2$-Wert fällt dagegen gleichmäßig weiter. Für $-3+0{,}3i$ ist der Effekt schwächer.
    Das Gewicht wirkt also, reicht aber an den Ecken nicht: Diese verlangen nach dem
    Vergleichskegel $\beta_{\mathrm{Ecke}}>0{,}39$ im \emph{Eckabstand}. Ein kombiniertes
    Kanten-Ecken-Gewicht ist der nächste Schritt.
  \item \emph{GMRES:} Für Glas konstant; für die plasmonischen Kontraste etwa $+6$
    Iterationen je Gradierungsstufe, wie im 2D-Prototyp ohne Eckblöcke. Erwartet ist, dass
    eine Kanten- und Eck-Blockvorkonditionierung (Schritt 2.5) dieses Wachstum beseitigt.
\end{itemize}

\paragraph{Grenzen.} $L\le5$ ($h_{\min}\approx0{,}03$) wegen 3\,GB Arbeitsspeicher. Die
Raten sind deshalb präasymptotisch, und die Aussage „beschränkt“ für Gold und für die
gewichteten Werte ist eine Tendenz, kein Nachweis. Die Anregung (ebene Welle) regt nicht
alle Sektoren an; die GMRES-Zahlen beziehen sich auf die angeregten.

\section{Kanten- und Eck-Block-Vorkonditionierung (Schritt 2.5)}\label{sec:vorkond3d}

\paragraph{Konstruktion.} Die Vorkonditionierung arbeitet je Sektor mit Block-Jacobi.
Im Fundamentaloktanten gibt es eine Ecke $(1,1,1)$ und drei Kantenstücke. Alle Dreiecke
mit Abstand $<R$ zur Ecke bilden einen Block. Die übrigen Dreiecke mit Abstand $<R$ zu
einer Kante bilden je Kante einen Block. Für alle anderen Dreiecke wird der eigene
$8\times8$-Diagonalblock invertiert. Die Blöcke werden exakt per LU invertiert, und
GMRES läuft rechtsvorkonditioniert (\texttt{cube\_precond.py}).

\paragraph{Iterationen} (bis $10^{-8}$, Maximum über die angeregten Sektoren):
\begin{center}\small
\begin{tabular}{llccccc}
\toprule
$\kappa$ & Vorkonditionierung & $L=1$ & 2 & 3 & 4 & 5\\
\midrule
Glas $2{,}25$ & ohne / $R=0{,}25$ / $R=0{,}5$ & 11/10/9 & -- & 11/10/9 & -- & 12/10/9\\
Gold $-11+1{,}2i$ & ohne & 43 & 48 & 50 & 54 & 56\\
 & nur Diagonalblöcke & 33 & 42 & 48 & 53 & 57\\
 & $R=0{,}25$ & 29 & 32 & 33 & 33 & 33\\
 & $R=0{,}5$ & 24 & 26 & 27 & 27 & 27\\
$-4+0{,}3i$ & ohne & 42 & 49 & 55 & 62 & 67\\
 & $R=0{,}25$ & 30 & 34 & 36 & 36 & 36\\
 & $R=0{,}5$ & 25 & 27 & 27 & 27 & 27\\
$-3+0{,}3i$ & ohne & 41 & 50 & 58 & 65 & 72\\
 & nur Diagonalblöcke & 35 & 46 & 56 & 65 & 74\\
 & $R=0{,}25$ & 30 & 35 & 36 & 36 & 36\\
 & $R=0{,}5$ & 25 & 27 & 28 & 28 & 28\\
\bottomrule
\end{tabular}
\end{center}

Mit Kanten- und Eckblöcken sind die Iterationszahlen ab $L=3$ konstant, für alle
Kontraste. Das gilt auch innerhalb der $L^2$-kritischen Kurve ($-3+0{,}3i$, $-4+0{,}3i$),
wo $\sigma_{\min}$ in $L^2$ gegen null geht (Abschnitt~\ref{sec:wuerfel}). Die
punktweisen Diagonalblöcke allein helfen dagegen nicht, und das Wachstum bleibt. Wie
im 2D-Prototyp ist das Wachstum also ein lokaler Kanten- und Eckeffekt.

\paragraph{Wahl des Radius} ($\kappa=-3+0{,}3i$):
\begin{center}\small
\begin{tabular}{lccc}
\toprule
$R$ & $L=3$ & 4 & 5\\
\midrule
$0{,}0625$ & 50 & 54 & 55\\
$0{,}125$ & 45 & 45 & 45\\
$0{,}25$ & 36 & 36 & 36\\
$0{,}5$ & 28 & 28 & 28\\
\bottomrule
\end{tabular}
\end{center}
Nötig ist ein \emph{fester physikalischer} Radius, hier ab etwa $R=2^{-3}$ der
Kantenlänge. Er darf nicht mit den Elementgrößen schrumpfen; bei $R=0{,}0625$ steigen
die Zahlen noch leicht. Größere Radien senken die absolute Iterationszahl.

\paragraph{Kosten.} Weil das Tensornetz auch entlang der Kanten verfeinert, wachsen die
Blöcke mit $L$. Für $R=0{,}25$ hat der größte Block 32, 64, 144, 384 und 720
Unbekannte; bei $L=5$ sind das $42\,\%$ eines Sektors. Für realistische Netze sind
deshalb anisotrope Kantennetze nötig, die entlang der Kante nicht verfeinern. Die
Kantenblöcke sollten dann über die $\mathcal{H}$-LU aus AP~3.4 invertiert werden. Bei
großen Streuern ist der Anteil der Blöcke an allen Unbekannten klein, weil er mit der
Kantenlänge und nicht mit der Fläche wächst.

\paragraph{Folgerung für H1.} In der Form „mit Kanten- und Eck-Blockvorkonditionierung“
ist H1 am Würfel für alle getesteten Kontraste numerisch gestützt: Die Iterationszahlen
hängen nicht von der Eckauflösung ab. Das gilt für die Iterationszahlen; die
Genauigkeit im Resonanzfenster (AP~1) ist davon unberührt.

\section{Grenzen und nächste Schritte}
\begin{itemize}[leftmargin=*]
  \item \emph{Speicher:} Mit 3\,GB Arbeitsspeicher sind dichte Matrizen bis etwa
    $N=500$ (4000 Unbekannte) möglich. Größere Studien brauchen die Kompression aus AP~3
    oder den Rust-Kern (Schritt 2.7).
  \item \emph{Quadratur:} Die Außenintegration der Nahpaare mit fester Unterteilung
    konvergiert an gemeinsamen Kanten nur langsam (logarithmische Singularität).
    Für höhere Genauigkeit ist Sauter-Schwab-Quadratur vorzusehen.
  \item \emph{Schritt 2.3, zweiter Teil:} Die analytische Frage nach einer
    Gårding-Ungleichung für $T_1$, die über die hier beobachtete Stabilität hinaus
    Galerkin-Konvergenz garantieren würde.
  \item \emph{Schritt 2.4, Fortsetzung:} kombiniertes Kanten-Ecken-Gewicht, Gradierung
    auch zu den Ecken hin, und die Mellin-Theorie der echten Würfelecke (AP~1), um die
    beobachtete Rate für $-4+0{,}3i$ zu erklären.
  \item \emph{Schritt 2.5, Fortsetzung:} Anisotrope, nichtkonforme Kantennetze sind
    umgesetzt und mit der Kompression aus AP~3 kombiniert (AP-3-Arbeitspapier,
    Abschnitt 5): Die Iterationszahlen bleiben konstant (47 bis 50). Offen ist die
    $\mathcal{H}$-LU für große Blöcke.
\end{itemize}

\end{document}
